\documentclass[11pt]{article}
\usepackage{amsmath,amssymb,amsthm,geometry}\geometry{margin=1in}
\newtheorem{proposition}{Proposition}
\title{A core-local random activation deployment with correlated entries}
\date{17 November 2026}\begin{document}\maketitle
\paragraph{Scope.}
This is a concrete instance test for the frozen-public-driver/core/private-record restriction. It is not a proof that every owner-aware receiver satisfies that restriction or that the resulting hom categories close. It supplies actual raw states, law admission, request ownership, context coefficients and observation fields, and isolates the capabilities excluded by the restriction. The full shared joint remains retained metadata. Neither its records nor its law label become locally readable merely because they were forwarded.

\section{The fixed typed scope}
Use a singleton physical world, retaining the paid identity-world event as a genuine atom. Fix a run-qualified event $e$, two distinct globally old actors $a_1,a_2$, and a shared ticket selecting one prescribed paid command at each actor. Both commands are selected before sampling. The frozen public driver $E$ contains this common boundary, ticket, actor/sponsor identities and declared primitive types; its raw carrier is a singleton. It has no fresh private raw value and no global joint-law selector. The common ticket contains both active sponsors, but each leaf's request is owned by its own $a_i$.

Actor $i$ has paid carrier $x_i\in\{0,1,2\}$ with law $p=(1/2,1/2,0)$. At every paid atom, including $x_i=2$, it creates one entry request at the actual postworld. The request's existence and type are constant, its predecessor is $D_i=x_i$, and its sponsor is $a_i$. Its complete initializer carrier is $z_i\in\{0,1,2\}$ with family consisting of $p$ at every predecessor. If a separate all-zero conditional plan is part of the interface, prescribe this same $p$ at the zero predecessor; it is extra supplied data, not recovered by disintegration.

Allocate the new child by the primitive name $(e,a_i,\mathrm{child},\tau)$, with an additional typed actorless material occurrence for the initializer coefficient. The names are the same at all raw atoms and in every content grouping; none uses a content wrapper. Old private records keep tags $(e,a_i,\mathrm{paid})$ and $(e,a_i,\mathrm{entry})$. A later composite may use their union as its own local core, but does not convert another leaf's tagged records into public driver records.

The selected state at every raw pair is
\[
 s_i'(x_i,z_i)=(\text{old actor done with record }x_i,
                 \text{new child complete with output }z_i).
\]
The child exposes its single output at entry. Its prescribed ENTIRE terminal collection is $I$ with unit content coefficient, including the full unused-grade convention. Its initializer CONTEXT coefficient is independently nonunit below. The old actor's output was empty, so the new port is an actual paid-sponsored activation, not fabricated port-invariant passivity. No newborn paid step occurs. The global selected state and ordered output list are the tuple of these leaf states and $(z_1,z_2)$.

\section{Admitted whole joints and a correlated member}
The full carrier is all $81$ tuples $(x_1,x_2,z_1,z_2)\in\{0,1,2\}^4$. Local admission fixes each complete pair marginal to
\[
                  Q_{x_i,z_i}=p\otimes p.              \tag{1}
\]
This retains the paid marginal and the entire initializer conditional test. No extra hidden whole-block constraint is assumed. Admit ALL global laws satisfying (1); a product law proves nonemptiness. The family is convex, but its chosen members remain complete command-law labels, not a randomized scheduler.

A nonproduct admitted member is
\[
 Q(x_1,x_2,z_1,z_2)=\begin{cases}
 1/4,&x_1,x_2\in\{0,1\},\quad z_1=x_2,\ z_2=x_1,\\
 0,&\text{otherwise}.
 \end{cases}                                           \tag{2}
\]
Each $(x_i,z_i)$ marginal is product fair on four points, with the five other local atoms retained at weight zero. Thus (1) holds, although each entry seed equals the other actor's paid outcome almost surely. This is permitted shared-scope coupling, not sequential fresh sampling conditioned on the entire global paid prefix.

\begin{proposition}[Nontrivial core-local deployment]
This scope has two genuinely active actors, random owned entries, cross-actor joint correlations and retained zero atoms. All local request, admission and selected-state functions factor through the frozen driver and their own tagged paid/entry core. Their complete local law labels are (1). They never require read access to a sibling's fresh coordinate or to an unrelated reserve-law label.
\end{proposition}
\begin{proof}
Existence, type, world and sponsor are fixed by the ticket and $a_i$; the predecessor is $x_i$. The initializer tests involve only the local pair marginal (1). The displayed selected-state function uses exactly $(x_i,z_i)$ on ALL nine local raw pairs, independently of the values at the other leaf. Taking all compatible global extensions preserves these equations without making a local function read the remote coordinates. Formula (2) satisfies (1) by summing over the other fair bit. The correlation is therefore in the chosen global law, not a covert remote argument of the local successor function.\end{proof}

\section{Target deployment versus the successive-entry source}
The preceding family is the LOCAL-PAIR target regime: each child initializer is tested given its own producer record $x_i$. It is not the successive-entry regime in which the entire producer tuple $X=(x_1,x_2)$ is filled before the child entries. That source imposes
\[
             Q(z_i=\cdot\mid x_1,x_2)=p
\]
at every positive full-$X$ fibre, together with its declared all-raw fields. Every such source law satisfies (1) by conditional averaging, but (2) does not: at each positive $X$ fibre, $z_1$ and $z_2$ are the fixed opposite coordinates, not fair bits. Thus (2) is an admitted core-compatible TARGET deployment outside that successive-entry source. It witnesses strictness of the directed conditional-coarsening face, not equality of the two families. The entire $Q$ and its producer marginal remain retained plan metadata on both sides; only local template consumption is restricted. No zero conditional table is inferred by the averaging argument.

\section{Actual nonunit material and full content}
In augmented finite-dimensional real vector spaces, put $K_i(x)=\mathbb R^2$ with first-coordinate augmentation at every paid atom. Put $L_i(z)=\mathbb R^2$ for $z=0,1$ and $L_i(2)=\mathbb R^3$, with first-coordinate augmentation. Each complete entry factor occurs once INSIDE its producing owner's tensor
\[
 W_i(x_i,z_i)=K_i(x_i)\otimes L_i(z_i).
\]
The actual shared scope context is the slice product
\[
             W_1(x_1,z_1)\times_{\mathbb R}W_2(x_2,z_2).
\]
Its dimensions on retained raw atoms are $7,9,11$, respectively when zero, one or two of the entry coordinates equal $2$. All augmentations are surjective, so these are actual fibre-product dimensions. Zero law does not erase these coefficients or their maps. Copy/group readdressing retains the two owner factors; its actual legs are their projections followed by canonical within-owner tensor reassociation. No tensor/limit interchange or coefficient duplication is invoked.

For a binary outer content shape, the terminal full grade is $C\circ I\cong C$, with actual unit content maps independent of $K_i,L_i$. No conclusion about content coefficients is drawn from the augmented context dimensions. Every unused outer shape and its coefficient map remains in the separately supplied full $C$.

\section{What this deployment does and does not consume}
A downstream local adapter may deterministically observe/store its own exported pair $(x_i,z_i)$, and may create further requests whose typed predecessor is that pair. Upon context composition this pair is an internal core of that same original leaf; the frozen public driver stays $E$. Sibling pairs stay forwarded reserve. Rebracketing local adapters changes none of those original primitive tags. This is a concrete supply of the core fields needed by a prospective composition theorem, not a proof of that theorem for every adapter.

The example is an instance of the owner-transparent random activation doctrine with complete initializer blocks retained. Its additional core restriction genuinely EXCLUDES the following previously possible behaviors. A local singleton-row response $s_i'=x_{3-i}$ does not descend through $(E,x_i,z_i)$ on the full carrier: fix that core and vary $x_{3-i}$. This fails even if an admitted law happens to correlate $z_i=x_{3-i}$; almost-sure equality does not provide the missing all-raw square. A request whose presence/type/disclosure needs a sibling's current private outcome is likewise excluded unless that dependency is already in the declared public driver, which cannot be enlarged after drawing the scope. A prescope command selected by the hidden sibling/global-law label rather than the local core-law label is excluded even though it is selected before sampling. Controllers selecting later commands from an entire retained global history are not automatically core-local either. Complete global histories may remain stored as reserve, but their use for local control requires a separate compatible controller interface.

Shared old ticket facts and public sponsors are not excluded. Nor are whole-law correlations such as (2). What is prohibited is promoting forwarded private records or full global law labels into newly public arguments during polynomial or context rebracketing. Same-scope owner-aware activation without that prohibition is a strictly broader class.

\paragraph{Finite diagnostic.}
\texttt{checks/check\_core\_activation\_instance.py} enumerates all 81 raw assignments, verifies normalization and both complete local marginals, checks all-raw successor factorization, confirms the remote-read obstruction, and checks the three nonunit context dimensions. This supports the stated finite instance; categorical closure remains a separate proof obligation.
\end{document}
